% GENERATED FILE. Edit canonical lessons, then run npm run generate:books.
% canonical-source-hash: fnv1a64:828c59e856b210a4
% canonical-lessons: BN-C159-tai, BN-C159-jehetu, BN-C159-jodi, BN-C159-tobu, BN-C159-nahole

\chapter{So, Since, If, Yet, Otherwise}
\label{ch:bn-so-since-if-yet-otherwise}
\hlchaptermodality{159}

\begin{chapteropening}
\textbf{By the end of this chapter:} \emph{I can say so, since, if, yet and otherwise.}

Complete the last lesson of chapter 159: I can say so, since, if, yet and otherwise.
\end{chapteropening}

\section[\texorpdfstring{\bn{তাই} (tāi)}{তাই (tāi)}]{\bn{তাই} (tāi) — so, therefore}

\label{lesson:BN-C159-tai}



\begin{quote}
Before the new one: say the Bengali for a minaret, then the Bengali for a ceiling.
\end{quote}

\subsection*{You'll want to know: \bn{তাই}}
\textbf{\bn{তাই}} — \emph{tāi} — \textquotedblleft{}so, therefore\textquotedblright{}.

\begin{cousinweb}
\textbf{\bn{তাই}} (\emph{tāi}) is \emph{tā}, \textquotedblleft{}that\textquotedblright{}, + \emph{-i}, \textquotedblleft{}indeed\textquotedblright{}: \textquotedblleft{}that very thing, so\textquotedblright{}. CDIAL lists Bengali \emph{tā} under the Sanskrit pronoun base \emph{ta-}, \textquotedblleft{}that\textquotedblright{} (CDIAL 5612), with Middle Bengali \emph{tẽ}, \textquotedblleft{}therefore\textquotedblright{}.
\end{cousinweb}

\begin{grammarlens}[title={What you've built}]
One more word for this chapter.
\end{grammarlens}

\subsection*{Your turn}
\begin{itemize}
\raggedright
  \item \emph{Say it:} \emph{tāi}
  \item \emph{Say it:} \emph{tāi}, once more
  \item \emph{Say it:} say \emph{tāi}
  \item \emph{Recall:} say the Bengali for a minaret, then the Bengali for a ceiling, then say \emph{tāi} again
\end{itemize}

\begin{tcolorbox}[breakable,colback=teal!4,colframe=teal!35!black,title={Before you move on}]
What does \bn{তাই} mean? (\textquotedblleft{}So, therefore\textquotedblright{}.) Say it once more.
\end{tcolorbox}
\section[\texorpdfstring{\bn{যেহেতু} (jehetu)}{যেহেতু (jehetu)}]{\bn{যেহেতু} (jehetu) — since, as}

\label{lesson:BN-C159-jehetu}



\begin{quote}
Before the new one: say the Bengali for a ceiling, then the Bengali for so.
\end{quote}

\subsection*{You'll want to know: \bn{যেহেতু}}
\textbf{\bn{যেহেতু}} — \emph{jehetu} — \textquotedblleft{}since, as\textquotedblright{}.

\begin{cousinweb}
\textbf{\bn{যেহেতু}} (\emph{jehetu}) is \emph{je}, \textquotedblleft{}which\textquotedblright{}, from the Sanskrit pronoun \emph{ya-}, \textquotedblleft{}who, which\textquotedblright{} (CDIAL 10391), + \emph{hetu}, \textquotedblleft{}cause\textquotedblright{}, Sanskrit \emph{hetu} taken over whole: \textquotedblleft{}for which cause\textquotedblright{}.
\end{cousinweb}

\begin{grammarlens}[title={What you've built}]
One more word for this chapter.
\end{grammarlens}

\subsection*{Your turn}
\begin{itemize}
\raggedright
  \item \emph{Say it:} \emph{jehetu}
  \item \emph{Say it:} \emph{jehetu}, once more
  \item \emph{Say it:} say \emph{jehetu}
  \item \emph{Recall:} say the Bengali for a ceiling, then the Bengali for so, then say \emph{jehetu} again
\end{itemize}

\begin{tcolorbox}[breakable,colback=teal!4,colframe=teal!35!black,title={Before you move on}]
What does \bn{যেহেতু} mean? (\textquotedblleft{}Since, as\textquotedblright{}.) Say it once more.
\end{tcolorbox}
\section[\texorpdfstring{\bn{যদি} (jôdi)}{যদি (jôdi)}]{\bn{যদি} (jôdi) — if}

\label{lesson:BN-C159-jodi}



\begin{quote}
Before the new one: say the Bengali for so, then the Bengali for since.
\end{quote}

\subsection*{You'll want to know: \bn{যদি}}
\textbf{\bn{যদি}} — \emph{jôdi} — \textquotedblleft{}if\textquotedblright{}.

\begin{cousinweb}
\textbf{\bn{যদি}} (\emph{jôdi}) is Sanskrit \emph{yadi}, \textquotedblleft{}if\textquotedblright{}, taken over whole as a learned loan, a tatsama (CDIAL 10410). The word inherited from it is Old Bengali \emph{jaï}, \textquotedblleft{}when\textquotedblright{}, which CDIAL lists.
\end{cousinweb}

\begin{grammarlens}[title={What you've built}]
One more word for this chapter.
\end{grammarlens}

\subsection*{Your turn}
\begin{itemize}
\raggedright
  \item \emph{Say it:} \emph{jôdi}
  \item \emph{Say it:} \emph{jôdi}, once more
  \item \emph{Say it:} say \emph{jôdi}
  \item \emph{Recall:} say the Bengali for so, then the Bengali for since, then say \emph{jôdi} again
\end{itemize}

\begin{tcolorbox}[breakable,colback=teal!4,colframe=teal!35!black,title={Before you move on}]
What does \bn{যদি} mean? (\textquotedblleft{}If\textquotedblright{}.) Say it once more.
\end{tcolorbox}
\section[\texorpdfstring{\bn{তবু} (tôbu)}{তবু (tôbu)}]{\bn{তবু} (tôbu) — yet, even so}

\label{lesson:BN-C159-tobu}



\begin{quote}
Before the new one: say the Bengali for since, then the Bengali for if.
\end{quote}

\subsection*{You'll want to know: \bn{তবু}}
\textbf{\bn{তবু}} — \emph{tôbu} — \textquotedblleft{}yet, even so\textquotedblright{}.

\begin{cousinweb}
The origin of \textbf{\bn{তবু}} (\emph{tôbu}) is not settled. CDIAL lists Assamese \emph{teo}, \textquotedblleft{}yet, still\textquotedblright{}, under Sanskrit \emph{tathāpi}, \textquotedblleft{}even so\textquotedblright{} (CDIAL 5647), and Bengali \emph{tabe}, \textquotedblleft{}then\textquotedblright{}, under \emph{ēvam ēva}, \textquotedblleft{}exactly so\textquotedblright{} (CDIAL 2528), but gives no Bengali \emph{tôbu}.
\end{cousinweb}

\begin{grammarlens}[title={What you've built}]
One more word for this chapter.
\end{grammarlens}

\subsection*{Your turn}
\begin{itemize}
\raggedright
  \item \emph{Say it:} \emph{tôbu}
  \item \emph{Say it:} \emph{tôbu}, once more
  \item \emph{Say it:} say \emph{tôbu}
  \item \emph{Recall:} say the Bengali for since, then the Bengali for if, then say \emph{tôbu} again
\end{itemize}

\begin{tcolorbox}[breakable,colback=teal!4,colframe=teal!35!black,title={Before you move on}]
What does \bn{তবু} mean? (\textquotedblleft{}Yet, even so\textquotedblright{}.) Say it once more.
\end{tcolorbox}
\section[\texorpdfstring{\bn{নাহলে} (nāhôle)}{নাহলে (nāhôle)}]{\bn{নাহলে} (nāhôle) — otherwise}

\label{lesson:BN-C159-nahole}



\begin{quote}
Before the new one: say the Bengali for if, then the Bengali for yet.
\end{quote}

\subsection*{You'll want to know: \bn{নাহলে}}
\textbf{\bn{নাহলে}} — \emph{nāhôle} — \textquotedblleft{}otherwise\textquotedblright{}.

\begin{cousinweb}
\textbf{\bn{নাহলে}} (\emph{nāhôle}) is \emph{nā}, \textquotedblleft{}not\textquotedblright{}, from Sanskrit \emph{na} (CDIAL 6906), + \emph{hôle}, \textquotedblleft{}if it is\textquotedblright{}, a form of \emph{hôoyā}, \textquotedblleft{}to be\textquotedblright{}, from Sanskrit \emph{bhavati}, \textquotedblleft{}becomes, is\textquotedblright{} (CDIAL 9416): \textquotedblleft{}if not so\textquotedblright{}.
\end{cousinweb}

\begin{grammarlens}[title={What you've built}]
One more word for this chapter.
\end{grammarlens}

\subsection*{Your turn}
\begin{itemize}
\raggedright
  \item \emph{Say it:} \emph{nāhôle}
  \item \emph{Say it:} \emph{nāhôle}, once more
  \item \emph{Say it:} say \emph{nāhôle}
  \item \emph{Recall:} say the Bengali for if, then the Bengali for yet, then say \emph{nāhôle} again
\end{itemize}

\begin{tcolorbox}[breakable,colback=teal!4,colframe=teal!35!black,title={Before you move on}]
What does \bn{নাহলে} mean? (\textquotedblleft{}Otherwise\textquotedblright{}.) Say it once more.
\end{tcolorbox}
